\chapter{Testing and Evaluation}\label{ch:evaluation}

\section{Test strategy}
The tests follow the structure of the system. The rule engine is a set of pure functions, so every
check kind is tested with invented documents and a fixed trip, including the boundary cases (a
passport that expires exactly three months after the return, an insurance that ends one day
early). The requirement files are tested for consistency: every rule cites a known source, every
field it reads exists with the right type, and no rule reads a document that no checklist item
asks for. Integration tests run the API against a real PostgreSQL database inside a transaction
that is rolled back; they cover registration and login, ownership (another user's file answers
\emph{not found}), encryption at rest, deletion of files together with their rows, the check of an
application file and the demo files. The web client is tested with Vitest and Testing Library,
including a test that fails when the Turkish and English translation files differ in their keys.

On 20 November 2026 the backend has 216 tests and the web client 39. Every pull request runs them
in GitHub Actions, and the \texttt{main} branch accepts a change only when both checks pass. Two
lessons from the tests themselves: tests must not depend on the developer's settings file, and
they must count only the rows they create, because the development database also holds the demo
account. Both were found when tests passed in CI but failed on the development machine, and were
fixed in the tests, not in the code.

The flow was also tried by hand in Chrome and Safari on a laptop and on a phone, including a
photograph taken with the phone's camera; Section~\ref{sec:web} lists what this found.

\section{Fall evaluation}
The fall success criteria (Section~\ref{sec:success}) compare the prototype with a reference or an
alternative. This section reports what is measured on 20 November 2026 and how the remaining
comparisons are made.

\subsection{Rule engine on the demo files}\label{sec:demo-eval}
\paragraph{Method.} The two demo files of Section~\ref{sec:demo} are first checked as they are,
then once for every labelled error, with that error alone applied. An error counts as
\emph{detected} when exactly its labelled finding appears: the named checklist item is missing, or
the named rule fails with the named issue code. Any other finding in the same run counts as a
false alarm. The command \texttt{python -m app.demo.evaluate} runs this and prints the tables below;
the same check runs in CI.

\paragraph{Result.} The error-free files give no finding: no missing document and no failed or
incomplete rule, with 16 rules checked for the employed applicant to the Netherlands and 10 for the
student to Norway. Every one of the 17 labelled errors is detected, and none of them produces a
further finding (Table~\ref{tab:demo-errors}).

{\small
\begin{longtable}{@{}p{0.9cm}>{\raggedright\arraybackslash}p{6.3cm}>{\raggedright\arraybackslash}p{5.3cm}p{1.4cm}@{}}
  \caption{Labelled errors of the demo files and the outcome (E: employed, Netherlands;
    S: student, Norway).}\label{tab:demo-errors}\\
  \toprule
  Error & Change & Expected finding & Outcome \\
  \midrule
  \endfirsthead
  \toprule
  Error & Change & Expected finding & Outcome \\
  \midrule
  \endhead
  \bottomrule
  \endfoot
  E01 & The travel insurance is not uploaded. & missing \reqid{nl-\allowbreak{}req-\allowbreak{}travel-\allowbreak{}insurance} & detected \\
  E02 & The passport expires two months after the return, one month too early. & \reqid{nl-\allowbreak{}rule-\allowbreak{}passport-\allowbreak{}validity} & detected \\
  E03 & The photo was taken eight months before the application. & \reqid{nl-\allowbreak{}rule-\allowbreak{}photo-\allowbreak{}age} & detected \\
  E04 & The insurance covers EUR 20,000 instead of at least EUR 30,000. & \reqid{nl-\allowbreak{}rule-\allowbreak{}insurance-\allowbreak{}amount} & detected \\
  E05 & The insurance ends one day before the return. & \reqid{nl-\allowbreak{}rule-\allowbreak{}insurance-\allowbreak{}period} & detected \\
  E06 & The hotel booking is in another person's name. & \reqid{eu-\allowbreak{}rule-\allowbreak{}names-\allowbreak{}match} & detected \\
  E07 & The approved leave ends two days before the return. & \reqid{eu-\allowbreak{}rule-\allowbreak{}leave-\allowbreak{}covers-\allowbreak{}trip} & detected \\
  E08 & The payslip of January is missing, so the three months are not consecutive. & \reqid{eu-\allowbreak{}rule-\allowbreak{}payslips-\allowbreak{}last-\allowbreak{}3-\allowbreak{}months} & detected \\
  E09 & The bank statement covers only two months. & \reqid{eu-\allowbreak{}rule-\allowbreak{}bank-\allowbreak{}statement-\allowbreak{}period} & detected \\
  E10 & The passport number in the employer letter has a wrong digit. & \reqid{eu-\allowbreak{}rule-\allowbreak{}passport-\allowbreak{}number-\allowbreak{}matches} & detected \\
  S01 & The bank statement is in Turkish and has no translation. & \reqid{no-\allowbreak{}rule-\allowbreak{}translation} & detected \\
  S02 & The civil registry extract was issued seven months before the application. & \reqid{no-\allowbreak{}rule-\allowbreak{}document-\allowbreak{}age} & detected \\
  S03 & The QR code of the YÖK certificate cannot be read. & \reqid{eu-\allowbreak{}rule-\allowbreak{}yok-\allowbreak{}qr-\allowbreak{}readable} & detected \\
  S04 & The proof of accommodation is not uploaded. & missing \reqid{eu-\allowbreak{}req-\allowbreak{}accommodation} & detected \\
  S05 & The insurance is valid in Norway only, not in the whole Schengen area. & \reqid{no-\allowbreak{}rule-\allowbreak{}insurance-\allowbreak{}territory} & detected \\
  S06 & The passport has one empty page. & \reqid{no-\allowbreak{}rule-\allowbreak{}passport-\allowbreak{}blank-\allowbreak{}pages} & detected \\
  S07 & The flight is booked in another person's name. & \reqid{eu-\allowbreak{}rule-\allowbreak{}names-\allowbreak{}match} & detected \\
\end{longtable}
}

\paragraph{What the result shows and what it does not.} The result shows that the rules implement
what they were written to check, that the engine compares dates with the trip and the application
date as intended, and that one error does not set off other rules. It does not show that the rules
match what a consulate decides: the labels and the rules were written by the same person from the
same sources, so a rule that misreads a source would be matched by an equally wrong label. The
comparison with a checklist compiled independently from the sources
(Section~\ref{sec:checklist-eval}) and, in spring, with human reviewers address this. Seventeen
errors in two files are also too few for rates: the result is a functional check, not a measured
accuracy.

\subsection{Generated checklist against a hand-compiled checklist}\label{sec:checklist-eval}
For each pilot destination and profile, a checklist is compiled by hand from the official pages,
without looking at the requirement files, and compared item by item with the generated one:
items missing from either side, items that differ in how they are submitted, and items whose
citation does not support them. This comparison is made in the week of 30 November 2026.

\subsection{The same model with and without the rule engine}\label{sec:model-comparison}
Suggested by the advisor on 20 November 2026, this comparison is the fall version of AS1. Both
sides receive the same document images rendered from the demo files, with every error variant:
\begin{itemize}
  \item \textbf{Model alone.} The model receives all documents of a file with the trip and the
    destination and is asked which documents are missing and which disagree with each other or
    with the requirements.
  \item \textbf{Model with the rule engine.} The model reads the fields of each document, and the
    rule engine of Section~\ref{sec:engine} decides.
\end{itemize}
Both are scored as in Section~\ref{sec:demo-eval}: labelled errors found, false alarms on the
error-free files and the error variants, time and cost per file. Each run is repeated three times
because model answers vary, and the model version and date are reported. The comparison is made
in the week of 30 November 2026.

\subsection{A general chat assistant}\label{sec:chat-eval}
The same question about the documents for one pilot profile is put to two general chat assistants
(ChatGPT and Claude), with the date and the product version recorded. Each answer is compared with
the hand-compiled checklist: missing items, wrong items, items without a source, and items taken
from another destination, such as a translation requirement for the Netherlands or none for
Norway. This also provides the evidence the advisor asked for behind the claims of
\texttt{docs/positioning.md}.

\section{Threats to validity}
\begin{itemize}
  \item \textbf{Synthetic documents.} The demo files are invented and clean; real documents are
    scanned, photographed at an angle, partly unreadable or in another layout. Results on them may
    not transfer to real applications.
  \item \textbf{Entered values.} In the evaluation of Section~\ref{sec:demo-eval} the values are
    entered correctly by construction, so it measures the rules, not the reading of documents.
  \item \textbf{Assumptions where sources are silent.} The engine assumes that a bank statement
    ends at most one month before the application, that the newest payslip is from the
    application month or the two before, and that amounts in another currency are not compared;
    for Norway the application form, passport and photograph are exempt from translation
    (ADR 0006). A consulate may apply these differently.
  \item \textbf{One author.} Requirements, rules, labels and evaluation were prepared by the same
    person; the independent checklist and the human reviewers of the spring reduce this.
  \item \textbf{Changing assistants.} General assistants and hosted models change over time; a
    result holds for the version and date reported.
\end{itemize}

\section{Spring experiment design}
\begin{description}[style=nextline,leftmargin=1.2cm]
  \item[AS1] A generator produces about 200 labelled synthetic files with more than 1,000
    documents, rendered with scanning artefacts; a test set is frozen before any tuning. The
    hybrid system and the model alone are compared per error type with precision, recall, F1,
    false-alarm rate, time and cost per file.
  \item[AS2] Human reviewers check a subset of the test set; their accuracy and time are compared
    with the hybrid system's on the same files.
  \item[AS3] Pairs of snapshots of the official pages are labelled meaningful, minor or cosmetic;
    the change detector is compared with a plain text-diff baseline with precision, recall, F1,
    false-alarm rate and detection delay.
\end{description}
